\section*{Hoofdstuk \ref{ch:b2:chromatography} --- Chromatografie}

\begin{solution}{exo:b2:chromatography:1}
$k = (5.0 - 1.0)/1.0 = 4.0$; fractie van de tijd in de \omterm{def:b2:chromatography:principle}{mobiele fase} $1/(1 + k) = 0.20$.
\end{solution}

\begin{solution}{exo:b2:chromatography:2}
$N = 16(6.0/0.30)^2 = 16 \times 400 = 6400$.
\end{solution}

\begin{solution}{exo:b2:chromatography:3}
$H = \qty{250}{mm}/10\,000 = \qty{0.025}{mm} = \qty{25}{\micro m}$.
\end{solution}

\begin{solution}{exo:b2:chromatography:4}
Ethanol in bloed: GC (vluchtig). Eiwit: HPLC. Cafeïne in een drank: HPLC (in water, zonder voorbehandeling niet vluchtig
genoeg). Benzeen in benzine: GC. Suikers: HPLC (ze ontleden voor ze koken).
\end{solution}

\begin{solution}{exo:b2:chromatography:5}
$R_s = 2(4.6 - 4.0)/(0.40 + 0.50) = 1.2/0.90 \approx 1.3$: niet helemaal tot op de basislijn gescheiden (onder 1.5).
\end{solution}

\begin{solution}{exo:b2:chromatography:6}
$\sqrt N = 4R_s\,\dfrac{\alpha}{\alpha - 1}\,\dfrac{1 + k_2}{k_2} = 4 \times 1.5 \times 11 \times 1.25 =
82.5$, dus $N \approx 6.8 \times 10^3$.
\end{solution}

\begin{solution}{exo:b2:chromatography:7}
$u_{\mathrm{opt}} = \sqrt{8/2} = \qty{2}{mm/s}$; $H_{\min} = 5 + 2\sqrt{16} = \qty{13}{\micro m}$.
\end{solution}

\begin{solution}{exo:b2:chromatography:8}
Fenol (polaire \ce{OH}, waterstofbruggen met water) eerst, dan benzeen, dan tolueen (één \ce{CH3}
meer, apolairder). Meer methanol maakt de \omterm{def:b2:chromatography:principle}{mobiele fase} minder polair: alle drie elueren vroeger, in
dezelfde volgorde.
\end{solution}

\begin{solution}{exo:b2:chromatography:9}
$F = (860/400)/(20.0/10.0) = 2.15/2.00 = 1.075$. Monster: $c_X = (645/410) \times 10.0/1.075 \approx
\qty{14.6}{mg/L}$.
\end{solution}

\begin{solution}{exo:b2:chromatography:10}
$R_s \propto \sqrt N \propto \sqrt L$: $L$ moet vermenigvuldigd worden met $(1.5/1.06)^2 \approx 2.0$, wat de
analysetijd en de druk verdubbelt. Andere hefbomen: de selectiviteit (\omterm{def:b2:chromatography:principle}{mobiele fase}, \omterm{def:b2:chromatography:principle}{stationaire fase},
temperatuur), kleinere deeltjes, het optimale debiet.
\end{solution}

\begin{solution}{exo:b2:chromatography:11}
De pieken liggen $4\sigma R_s$ uit elkaar, dus het midden ligt op $2\sigma R_s$ van elk. $R_s = 1$: elke piek
draagt $\mathrm e^{-(2)^2/2} = \mathrm e^{-2} = 0.135$ bij; het dal ligt op $0.27$ van de piekhoogte.
$R_s = 1.5$: $2\mathrm e^{-(3)^2/2} = 2\mathrm e^{-4.5} \approx 0.022$, ongeveer 2~\%: in de praktijk terug op de basislijn.
\end{solution}

\begin{solution}{exo:b2:chromatography:12}
$k_2/(1+k_2)$: 0.67, 0.83, 0.91; tijd in eenheden $t_M$: 3, 6, 11. Van 2 naar 5 gaan wint 25~\% aan
\omterm{def:b2:chromatography:separation}{resolutie} voor twee keer de tijd; van 5 naar 10 9~\% voor opnieuw bijna twee keer de tijd. Een $k$ van ongeveer 2 tot 5 is
het gebruikelijke compromis.
\end{solution}

\begin{solution}{pb:b2:chromatography:1}
\textbf{1.} Stationair: silicakorrels met lange alkylketens, apolair. Mobiel: water met
methanol, polair.
\textbf{2.} Ze zijn zeer polair en blijven in de \omterm{def:b2:chromatography:principle}{mobiele fase}: $k \approx 0$.
\textbf{3.} Cafeïne heeft één methylgroep meer dan theofylline, en theofylline een \ce{N-H} die
waterstofbruggen vormt met water: theofylline is polairder en elueert eerst.
\textbf{4.} Ze lijkt qua structuur op cafeïne (vergelijkbare respons en vergelijkbaar gedrag), is afwezig in de drank, en elueert
dicht bij cafeïne terwijl ze ervan gescheiden blijft.
\textbf{5.} Beide zouden dalen: een minder polaire \omterm{def:b2:chromatography:principle}{mobiele fase} concurreert beter met de \omterm{def:b2:chromatography:principle}{stationaire fase}.
\textbf{6.} De tijd die de \omterm{def:b2:chromatography:principle}{mobiele fase} nodig heeft om de kolom te doorlopen, de tijd die een verbinding in beweging doorbrengt.
\textbf{7.} $k_{\mathrm{theo}} = (3.0 - 1.0)/1.0 = 2.0$; $k_{\mathrm{caf}} = 2.4$.
\textbf{8.} $\alpha = 2.4/2.0 = 1.2$.
\textbf{9.} Cafeïne $N = 16(3.4/0.20)^2 = 4624 \approx 4.6 \times 10^3$; theofylline $16(3.0/0.18)^2
\approx 4.4 \times 10^3$.
\textbf{10.} $H = \qty{150}{mm}/4624 \approx \qty{0.032}{mm} = \qty{32}{\micro m}$.
\textbf{11.} $R_s = 2(3.4 - 3.0)/(0.18 + 0.20) = 0.80/0.38 \approx 2.1$.
\textbf{12.} $\frac{\sqrt{4624}}{4} \times \frac{0.2}{1.2} \times \frac{2.4}{3.4} = 17 \times 0.167 \times
0.706 \approx 2.0$; het kleine verschil komt van de aanname van gelijke breedten in de vergelijking.
\textbf{13.} Ja: $R_s > 1.5$.
\textbf{14.} $N$ wordt gehalveerd: $R_s \approx 2.1/\sqrt2 \approx 1.5$, net basislijnscheiding.
\textbf{15.} Gehalveerd: cafeïne bij \qty{1.7}{min}.
\textbf{16.} $R_s \approx 2.1\sqrt2 \approx 3.0$, voor twee keer de tijd en twee keer de druk: hier verspilling.
\textbf{17.} De samenstelling van de \omterm{def:b2:chromatography:principle}{mobiele fase} (methanolfractie, pH, een ander organisch oplosmiddel) of de
\omterm{def:b2:chromatography:principle}{stationaire fase}.
\textbf{18.} Weinig: $k_2/(1 + k_2) = 0.71$ al; bij $k_2 = 5$ zou het 0.83 zijn (+18~\%) voor een
analysetijd die vermenigvuldigd wordt met $6/3.4 \approx 1.8$.
\textbf{19.} $F = (1500/1000)/(50.0/40.0) = 1.50/1.25 = 1.20$.
\textbf{20.} Beide pieken komen uit dezelfde injectie: het volume valt weg in de verhouding van de oppervlakten.
\textbf{21.} $c = (970/1010) \times 40.0/1.20 \approx \qty{32.0}{mg/L}$.
\textbf{22.} Tienvoudige verdunning: \qty{320}{mg/L}.
\textbf{23.} $A_{IS}$ zou te groot zijn, en het resultaat voor cafeïne te laag.
\textbf{24.} Een reeks cafeïnestandaarden, geïnjecteerd onder exact dezelfde omstandigheden als het monster, een
ijklijn van oppervlakte tegen concentratie, en reproduceerbare injectievolumes.
\textbf{25.} $\qty{320}{mg/L} \times \qty{0.250}{L}$: $\boldsymbol{\approx \qty{80}{mg}}$ cafeïne per blikje (gegevens
voor de oefening).
\end{solution}
